The planned direct quantization of the archived four-layer ONNX graph failed during ONNX shape inference with a 768-versus-2 dimension conflict. The prespecified attempt ended at conversion. In a separately labeled exploratory repair, I cleared the intermediate tensor annotations before quantization. This left the FP32 operators, weights, and input and output definitions intact. The original and stripped FP32 graphs agreed within \(10^{-4}\) absolute and relative tolerance on probe inputs of 3, 16, 32, 55, 211, and 512 tokens. The repaired input allowed ORT to construct the INT8 candidate with the original quantizer settings.

The student graph shrank from 211.6 MB to 124.8 MB, a 41.0\% reduction in serialized bytes. On SST-2 the INT8 student scored 760/872 (87.16\%), versus 765/872 (87.73\%) for FP32. Its frozen confidence cascade scored 785/872 (90.02\%), versus 789/872 (90.48\%) for FP32 and 794/872 (91.06\%) for the full teacher. The exact loss against the teacher was \(100(794-785)/872=1.032\) points, exceeding the one-point target. Quantization changed 27 student predictions and 61 routing decisions, raising teacher calls from 161 to 206. A smaller serialized graph alone cannot establish lower request latency or process memory.

On the reused IMDb sample, each student got 843/1,000 correct, with 40 predictions differing between FP32 and INT8. The planned resident timing run requires AC power and remains pending. The accuracy audit already rules out this INT8 cascade at the tested cutoff under the project's observed acceptance criteria.
